\documentclass[11pt,a4paper]{article}

% --- Packages ---
\usepackage[utf8]{inputenc}
\usepackage[T1]{fontenc}
\usepackage[english]{babel}
\usepackage{amsmath,amssymb,amsfonts}
\usepackage{graphicx}
\usepackage{booktabs}
\usepackage{hyperref}
\usepackage{cleveref}
\usepackage{xcolor}
\usepackage{float}
\usepackage[margin=2.5cm]{geometry}
\usepackage{authblk}
\usepackage{listings}
\usepackage{algorithm}
\usepackage{algpseudocode}
\usepackage{subcaption}
\usepackage{multirow}

\hypersetup{colorlinks=true,linkcolor=blue!60!black,citecolor=green!50!black,urlcolor=blue!60!black}

\lstset{
  basicstyle=\ttfamily\footnotesize,
  breaklines=true,
  frame=single,
  language=Python,
}

% --- Title ---
\title{Steering Vectors vs.\ Prompt Engineering for Agentic Calendar Tasks:\\
A Comparative Study on Instruction-Tuned LLMs}

\author[1]{[Author Name]}
\affil[1]{[Affiliation]}

\date{\today}

\begin{document}
\maketitle

% ===================================================================
\begin{abstract}
We present a comparative study of \textbf{activation steering} versus \textbf{prompt engineering} for improving task-specific behavior in instruction-tuned language models, focusing on an agentic calendar scheduling use case. Using Qwen3-4B-Instruct and its base variant as primary models, we (1)~analyze how four different tokenizers fragment calendar-related inputs, (2)~extract contrastive steering vectors from residual stream activations using a SAE-inspired methodology, (3)~perform a comprehensive layer-by-layer steering sweep across 11 layers and 6 coefficient levels, (4)~validate findings on the base (non-instruct) model, (5)~test budget guidance for reasoning length control, and (6)~map SAE features via Neuronpedia's transcoder dictionary.

Our key findings are: (a)~\textbf{mid-layer steering (layers 15--18) achieves 100\% behavioral change} at moderate coefficients ($\alpha=30$), while late layers (33--35) are rigid; (b)~\textbf{base models are more steerable but fragile}---the sweet spot at layer 15, $\alpha=30$ boosts calendar scores from 0.70 to 0.97, but higher coefficients cause degeneration; (c)~\textbf{budget guidance has no effect} on instruct models that already produce compact structured outputs ($\sim$75 tokens); (d)~Neuronpedia transcoder features confirm calendar-related circuits distribute across layers 20--35 with extreme sparsity (9/450 sampled features, 0.07\% density). We discuss implications for the choice between steering and prompting in agentic applications.
\end{abstract}

\textbf{Keywords:} activation steering, prompt engineering, tokenizer analysis, mechanistic interpretability, agentic LLM, calendar scheduling

% ===================================================================
\section{Introduction}
\label{sec:intro}

Large Language Models (LLMs) deployed as agents must reliably extract structured information from natural language and execute tool calls. A common real-world task is calendar event creation: given a user utterance like \textit{``Schedule a meeting tomorrow at 2pm with Marie''}, the agent must extract date, time, title, attendees, and location, then invoke a \texttt{create\_calendar\_event} function.

Two families of techniques can steer models toward task-specific behavior:

\begin{enumerate}
    \item \textbf{Prompt engineering} --- crafting system prompts, few-shot examples, and chain-of-thought instructions to guide generation at inference time.
    \item \textbf{Activation steering} --- adding learned or computed vectors to the model's residual stream during forward passes, modifying internal representations without changing weights~\cite{turner2023activation,li2024inference}.
\end{enumerate}

Recent work has shown that per-layer steering vectors trained via reinforcement learning can match fully fine-tuned models on mathematical reasoning tasks, adding only 0.0016\% parameters~\cite{gao2025steering}. Separately, budget guidance techniques demonstrate training-free control of reasoning length via token-level distribution modeling~\cite{li2025budget}. However, these approaches have primarily been studied on reasoning benchmarks --- their effectiveness in agentic tool-use contexts remains unexplored.

\textbf{Contributions.} This paper makes six contributions:
\begin{itemize}
    \item A comparative tokenizer analysis across four models revealing universal fragmentation of temporal expressions (\cref{sec:tokenizers}).
    \item SAE-inspired contrastive feature extraction with logit lens validation (\cref{sec:features}).
    \item A comprehensive steering sweep across 11 layers and 6 coefficients, identifying mid-layer steering (layers 15--18) as the effective intervention point (\cref{sec:steering}).
    \item Comparative steering on the base (non-instruct) model, revealing higher malleability but greater fragility (\cref{sec:base_steering}).
    \item Evidence that budget guidance is ineffective for already-compact structured output tasks (\cref{sec:budget}).
    \item Neuronpedia SAE feature mapping confirming extreme sparsity of calendar circuits (\cref{sec:neuronpedia}).
\end{itemize}

% ===================================================================
\section{Related Work}
\label{sec:related}

\subsection{Activation Steering}
Turner et al.~\cite{turner2023activation} introduced activation addition, showing that adding contrastive mean-difference vectors to residual streams can control model behavior (e.g., shifting sentiment or honesty). Gao et al.~\cite{gao2025steering} extended this to reinforcement learning, training per-layer bias vectors that match full RL fine-tuning on math reasoning with extreme parameter efficiency. Li et al.~\cite{li2025budget} proposed budget guidance, a training-free method using Gamma-distribution predictors to control reasoning length at inference time.

\subsection{Mechanistic Interpretability}
Sparse Autoencoders (SAEs) decompose model activations into interpretable features~\cite{cunningham2023sparse}. The Gemma Scope project~\cite{lieberum2024gemma} provides SAEs, transcoders, and crosscoders trained on every layer of Gemma models. Neuronpedia\footnote{\url{https://www.neuronpedia.org}} offers circuit-level attribution graphs for models including Qwen3-4B.

\subsection{Prompt Engineering for Tool Use}
Function calling and structured output extraction have been studied extensively in the context of LLM agents~\cite{schick2023toolformer}. Chain-of-thought prompting~\cite{wei2022chain} and few-shot examples remain strong baselines for structured extraction tasks.

% ===================================================================
\section{Tokenizer Analysis}
\label{sec:tokenizers}

Tokenization quality directly impacts downstream model behavior: a tokenizer that fragments semantic units distributes information across multiple positions in the residual stream, potentially diluting features that steering vectors target.

\subsection{Experimental Setup}

We compare four tokenizers on nine prompt categories designed for agentic calendar tasks:

\begin{table}[H]
\centering
\caption{Tokenizer properties.}
\label{tab:tokenizers}
\begin{tabular}{lcc}
\toprule
\textbf{Model} & \textbf{Vocab Size} & \textbf{Architecture} \\
\midrule
Qwen3-4B-Instruct-2507 & 151,643 & BPE \\
Gemma-3-1B-IT & 262,144 & SentencePiece \\
Llama-3.2-3B-Instruct & 128,000 & BPE \\
Phi-3-mini-4k-instruct & 32,000 & SentencePiece \\
\bottomrule
\end{tabular}
\end{table}

Our test suite includes simple and complex calendar requests in French and English, relative and absolute temporal expressions, JSON tool calls, and system prompts.

\subsection{Results}

\subsubsection{Token Counts and Compression}

\begin{table}[H]
\centering
\caption{Token counts across prompt categories. Lower is better.}
\label{tab:token_counts}
\small
\begin{tabular}{lcccc}
\toprule
\textbf{Prompt} & \textbf{Qwen3} & \textbf{Gemma-3} & \textbf{Llama-3.2} & \textbf{Phi-3} \\
\midrule
simple\_fr & 20 & 19 & 19 & 22 \\
simple\_en & 15 & 15 & 15 & 17 \\
complex\_fr & 55 & 50 & 51 & 62 \\
complex\_en & 47 & 45 & 43 & 53 \\
ambiguous & 18 & 16 & 18 & 22 \\
temporal\_relative & 23 & 20 & 24 & 27 \\
temporal\_absolute & 43 & 43 & 29 & 44 \\
tool\_call\_json & 125 & 125 & 106 & 136 \\
system\_agent & 34 & 36 & 34 & 41 \\
\bottomrule
\end{tabular}
\end{table}

\textbf{Llama-3.2-3B achieves the best compression} on both temporal expressions (29 tokens vs.\ 43--44 for ISO 8601 dates) and JSON tool calls (106 vs.\ 125--136). This is attributable to its BPE vocabulary including multi-digit numbers (\texttt{``24''}, \texttt{``202''}) and code-oriented subwords (\texttt{``\_calendar\_event''} as a single token).

\subsubsection{Quality Scoring}

We designed a composite scoring function evaluating four dimensions relevant to agentic calendar use:

\begin{table}[H]
\centering
\caption{Tokenizer quality scores for calendar tasks (0--1 scale, higher is better).}
\label{tab:quality_scores}
\begin{tabular}{lccccc}
\toprule
\textbf{Model} & \textbf{Temporal} & \textbf{Semantic} & \textbf{JSON} & \textbf{Multilingual} & \textbf{Overall} \\
\midrule
Qwen3-4B & 0.205 & 0.680 & 0.288 & 0.802 & 0.483 \\
Gemma-3-1B & 0.226 & 0.722 & 0.288 & 0.845 & 0.511 \\
Llama-3.2-3B & 0.316 & 0.726 & 0.375 & 0.816 & \textbf{0.551} \\
Phi-3-mini & 0.173 & 0.580 & 0.249 & 0.814 & 0.438 \\
\bottomrule
\end{tabular}
\end{table}

\textbf{Key finding:} All tokenizers score poorly on temporal integrity ($<0.32$). ISO 8601 timestamps are universally fragmented into individual characters, creating a distributed representation that complicates feature extraction for steering.

% ===================================================================
\section{Feature Extraction via Contrastive Activations}
\label{sec:features}

\subsection{Methodology}

Inspired by the Gemma Scope 2 pipeline~\cite{lieberum2024gemma} and contrastive activation addition~\cite{turner2023activation}, we extract steering-relevant features without training SAEs. Our approach:

\begin{enumerate}
    \item Collect last-token residual stream activations at all 36 layers for 10 calendar and 10 neutral prompts.
    \item Compute per-layer mean activations for each group.
    \item Define the \textbf{steering vector} as the mean difference: $\mathbf{v}_l = \bar{\mathbf{a}}_l^{\text{cal}} - \bar{\mathbf{a}}_l^{\text{neu}}$.
    \item Rank layers by $\|\mathbf{v}_l\|_2$ (separation magnitude) and cosine distance.
\end{enumerate}

\subsection{Layer Importance}

\begin{table}[H]
\centering
\caption{Top-10 layers by steering vector L2 norm.}
\label{tab:layer_importance}
\begin{tabular}{cccc}
\toprule
\textbf{Rank} & \textbf{Layer} & \textbf{L2 Norm} & \textbf{Cosine Distance} \\
\midrule
1 & 35 & 361.00 & 0.151 \\
2 & 34 & 218.63 & 0.075 \\
3 & 33 & 179.38 & 0.094 \\
4 & 32 & 162.50 & 0.104 \\
5 & 31 & 143.63 & 0.113 \\
6 & 30 & 126.19 & 0.120 \\
7 & 29 & 110.31 & 0.130 \\
8 & 28 & 100.44 & 0.151 \\
9 & 27 & 89.50 & 0.164 \\
10 & 26 & 79.50 & 0.168 \\
\bottomrule
\end{tabular}
\end{table}

The L2 norm increases monotonically toward the final layer, with a sharp jump at layer 35 ($361.0$ vs.\ $218.6$ at layer 34). This is consistent with the ``representation engineering'' hypothesis: later layers encode increasingly task-specific, high-level semantic features.

\subsection{Logit Lens Analysis}

We project the top-3 steering vectors through the unembedding matrix ($\mathbf{W}_U$) after applying RMSNorm weights:

$$\text{logits}_l = \mathbf{W}_U \cdot (\mathbf{w}_{\text{norm}} \odot \mathbf{v}_l)$$

\begin{table}[H]
\centering
\caption{Top tokens promoted/suppressed by steering vectors (layers 33--35).}
\label{tab:logit_lens}
\small
\begin{tabular}{cl}
\toprule
\textbf{Layer} & \textbf{Top Promoted Tokens} \\
\midrule
35 & \texttt{Schedule}, \texttt{Agenda}, \texttt{agenda}, \texttt{schedule}, \texttt{Scheduled} \\
34 & \texttt{schedule}, \texttt{agenda}, \texttt{attendees}, \texttt{RSVP}, \texttt{calendar} \\
33 & \texttt{agenda}, \texttt{schedule}, \texttt{attendees}, \texttt{invite}, \texttt{scheduled} \\
\midrule
\textbf{Layer} & \textbf{Top Suppressed Tokens} \\
\midrule
35 & \texttt{Magnitude}, \texttt{hesion}, \texttt{magnitude}, \texttt{density}, \texttt{humidity} \\
34 & \texttt{Magnitude}, \texttt{licking}, \texttt{entropy}, \texttt{Measured} \\
33 & \texttt{entropy}, \texttt{Measured}, \texttt{Magnitude} \\
\bottomrule
\end{tabular}
\end{table}

\textbf{Cross-lingual signal:} The steering vectors also promote Chinese tokens for ``email'' and ``modify'', indicating that Qwen3-4B has learned cross-lingual calendar-related associations. This is notable given the model's multilingual training on 119 languages.

% ===================================================================
\section{Steering Vector Application}
\label{sec:steering}

\subsection{Methodology}

We apply steering vectors using activation addition~\cite{turner2023activation}:

$$\mathbf{h}_l' = \mathbf{h}_l + \alpha \cdot \frac{\mathbf{v}_l}{\|\mathbf{v}_l\|}$$

where $\alpha \in \{0, 10, 30, 60, 100, 200\}$ and $\mathbf{h}_l$ is the residual stream at layer $l$. We sweep across 11 layers ($l \in \{5, 10, 15, 18, 20, 22, 25, 28, 30, 33, 35\}$) on six prompts: 2 calendar (FR/EN), 2 ambiguous (FR/EN), and 2 non-calendar (FR/EN).

\subsection{Initial Finding: Late-Layer Rigidity}

Applying the layer-35 vector (highest L2 norm, 361.0) produces \textbf{zero behavioral change} across all coefficients up to $\alpha=50$. All responses are character-for-character identical to baseline. This motivates the broader layer sweep.

\subsection{Mid-Layer Steering Sweet Spot}

\begin{table}[H]
\centering
\caption{Change rate (\%) vs.\ baseline across layers and coefficients on Qwen3-4B-Instruct. Bold indicates 100\% change rate with coherent outputs.}
\label{tab:midlayer_sweep}
\small
\begin{tabular}{rcccccc}
\toprule
& \multicolumn{6}{c}{\textbf{Steering Coefficient ($\alpha$)}} \\
\cmidrule(lr){2-7}
\textbf{Layer} & 0 & 10 & 30 & 60 & 100 & 200 \\
\midrule
5  & 0 & 33 & 33 & 100$^\dagger$ & 67$^\dagger$ & 67$^\dagger$ \\
10 & 0 & 17 & 67 & 83 & 67 & 100 \\
\textbf{15} & \textbf{0} & \textbf{17} & \textbf{100} & \textbf{100} & \textbf{83} & \textbf{100} \\
\textbf{18} & \textbf{0} & \textbf{33} & \textbf{100} & \textbf{100} & \textbf{100} & \textbf{67} \\
20 & 0 & 17 & 33 & 100 & 100 & 67 \\
22 & 0 & 17 & 17 & 50 & 100 & 67 \\
25 & 0 & 0 & 17 & 33 & 50 & 83 \\
30 & 0 & 0 & 17 & 17 & 17 & 50 \\
33 & 0 & 0 & 17 & 17 & 17 & 17 \\
35 & 0 & 0 & 0 & 0 & 0 & 17 \\
\bottomrule
\end{tabular}
\vspace{2pt}

{\footnotesize $^\dagger$ Layer 5 at high $\alpha$ produces degenerate (empty) outputs.}
\end{table}

\textbf{Key result:} Layers 15--18 at $\alpha=30$ achieve 100\% change rate with coherent, well-formed responses. This is the \textit{steering sweet spot}---representations at these layers are transitioning from syntactic to semantic and remain malleable enough for intervention.

\textbf{Qualitative effect:} Non-calendar prompts get reinterpreted as calendar tasks. For instance, \textit{``Explain how transformers work''} becomes \textit{``Let's schedule a meeting to discuss transformers''}---the steering vector genuinely redirects the model's intent, not just its surface tokens.

A clear \textbf{gradient of rigidity} emerges: late layers (30--35) require 3--10$\times$ higher coefficients for the same effect, and even then produce only partial changes. This confirms the hypothesis that instruction tuning ``freezes'' late-layer representations.

% ===================================================================
\section{Base Model Steering}
\label{sec:base_steering}

\subsection{Setup}

To disentangle the effect of instruction tuning, we repeat the steering experiment on \textbf{Qwen3-4B} (base, non-instruct). We extract fresh contrastive vectors using completion-style prompts (no chat template) with 10 calendar and 10 neutral examples.

\subsection{Results}

\begin{table}[H]
\centering
\caption{Calendar score (0--1) for base model steering. Higher = stronger calendar signal. Baseline avg = 0.70.}
\label{tab:base_steering}
\small
\begin{tabular}{rcccccc}
\toprule
& \multicolumn{6}{c}{\textbf{Steering Coefficient ($\alpha$)}} \\
\cmidrule(lr){2-7}
\textbf{Layer} & 0 & 10 & 30 & 60 & 100 & 200 \\
\midrule
10 & 0.70 & 0.73 & 0.70 & 0.73 & 0.23 & 0.00 \\
\textbf{15} & \textbf{0.70} & \textbf{0.70} & \textbf{0.97} & \textbf{0.33} & \textbf{0.03} & \textbf{0.00} \\
18 & 0.70 & 0.73 & 0.73 & 0.13 & 0.00 & 0.00 \\
20 & 0.70 & 0.73 & 0.73 & 0.03 & 0.00 & 0.00 \\
25 & 0.70 & 0.73 & 0.73 & 0.73 & 0.77 & 0.00 \\
30 & 0.70 & 0.70 & 0.70 & 0.73 & 0.73 & 0.70 \\
35 & 0.70 & 0.70 & 0.70 & 0.70 & 0.73 & 0.77 \\
\bottomrule
\end{tabular}
\end{table}

Three findings emerge:

\begin{enumerate}
    \item \textbf{Layer 15 at $\alpha=30$ is the optimal point}: calendar score jumps from 0.70 to 0.97 ($+0.27$), matching the instruct model's sweet spot layer.
    \item \textbf{Base models are much more fragile}: at $\alpha \geq 60$, scores collapse to near-zero (degenerate output). The instruct model tolerates $\alpha=200$ at mid-layers without complete degeneration.
    \item \textbf{Late-layer rigidity is architectural}: layers 30--35 resist steering on both instruct and base models, ruling out instruction tuning as the sole cause. This is a property of the transformer architecture itself---late residual streams are too committed to the output distribution.
\end{enumerate}

% ===================================================================
\section{Budget Guidance}
\label{sec:budget}

Inspired by Li et al.~\cite{li2025budget}, we implement a simplified budget guidance mechanism: a Gamma-distribution predictor that biases the EOS token logit as generation approaches a token budget, using a sigmoid urgency ramp.

\subsection{Results}

We test budgets of 32, 64, 128, 256, and 512 tokens on 5 calendar prompts:

\begin{table}[H]
\centering
\caption{Budget guidance results on Qwen3-4B-Instruct. Baseline avg: 75 tokens.}
\label{tab:budget}
\begin{tabular}{rcccc}
\toprule
\textbf{Budget} & \textbf{Avg Tokens} & \textbf{Valid} & \textbf{Within Budget} & \textbf{Savings} \\
\midrule
$\infty$ (baseline) & 74.6 & 100\% & 100\% & 0\% \\
32 & 74.6 & 100\% & 0\% & 0\% \\
64 & 74.6 & 100\% & 0\% & 0\% \\
128 & 74.6 & 100\% & 100\% & 0\% \\
256 & 74.6 & 100\% & 100\% & 0\% \\
512 & 74.6 & 100\% & 100\% & 0\% \\
\bottomrule
\end{tabular}
\end{table}

\textbf{Budget guidance has zero effect.} The instruct model already generates near-minimal outputs ($\sim$75 tokens of compact JSON). Even the most aggressive budget (32 tokens) fails to further compress the output---the model's greedy decoding path is already optimal for structured extraction. This contrasts with reasoning-heavy tasks (e.g., GSM8K) where models produce verbose chain-of-thought before answering, and where budget guidance can yield significant savings~\cite{li2025budget}.

\textbf{Implication:} Budget guidance targets \textit{thinking overhead}, not output verbosity. For tasks where instruction tuning has already eliminated unnecessary reasoning tokens, it adds no value.

% ===================================================================
\section{Neuronpedia Feature Mapping}
\label{sec:neuronpedia}

We explore Qwen3-4B's SAE features via Neuronpedia's transcoder dictionary (163,840 features per layer) using three strategies: explanation keyword search, random sampling, and targeted key-layer scans.

\subsection{Results}

\begin{itemize}
    \item \textbf{Explanation search}: 115 unique features matched calendar keywords (schedule, meeting, date, time, etc.) across layers 20--35.
    \item \textbf{Random sampling}: only 9 out of 450 sampled features (2\%) were calendar-related, confirming extreme sparsity.
    \item \textbf{Feature distribution}: calendar features concentrate in later layers (20--35), consistent with our contrastive analysis. The highest-layer features show the most specific explanations (e.g., ``scheduling appointments'' vs.\ generic ``time references'').
\end{itemize}

The 0.07\% density of calendar features (9/450 random samples $\times$ 163,840 features/layer) implies approximately 115 calendar-relevant features per layer---a tiny fraction of the total, but sufficient to encode the task-specific behavior observed in our steering experiments.

% ===================================================================
\section{Prompt Engineering Baselines}
\label{sec:baselines}

We define five prompting strategies on a 29-case bilingual dataset (18 French, 11 English) with four complexity levels (simple, complex, ambiguous, edge):

\begin{table}[H]
\centering
\caption{Prompt engineering strategies.}
\label{tab:strategies}
\begin{tabular}{lcc}
\toprule
\textbf{Strategy} & \textbf{Messages} & \textbf{Key Feature} \\
\midrule
Zero-shot & 2 & System prompt only \\
Few-shot (3) & 8 & 3 example pairs \\
Few-shot (5) & 12 & 5 example pairs \\
Chain-of-thought & 2 & Step-by-step reasoning \\
Tool use & 2 + tool & Function calling schema \\
\bottomrule
\end{tabular}
\end{table}

Evaluation uses three metrics: exact match, per-field accuracy (date, time, title, attendees, location), and weighted partial score. Full inference results with Ollama/llama.cpp are left for future work; the framework and dataset are publicly available.

% ===================================================================
\section{Discussion}
\label{sec:discussion}

\subsection{The Layer Hierarchy of Steerability}

Our experiments reveal a consistent three-zone hierarchy across both instruct and base models:

\begin{itemize}
    \item \textbf{Early layers (1--10):} Highly sensitive but unstable. Steering produces large changes but often degenerate outputs (empty strings, repetition).
    \item \textbf{Mid layers (15--18):} The \textit{sweet spot}. Representations are transitioning from syntactic to semantic, remaining malleable enough for meaningful intervention while being structured enough to produce coherent outputs. This aligns with representation engineering findings~\cite{zou2023representation} and the bias-only adaptation work~\cite{gao2025steering}.
    \item \textbf{Late layers (30--35):} Rigid. Despite having the highest L2 norms (i.e., the most discriminative vectors), these layers resist perturbation. We show this is \textit{architectural}, not a consequence of instruction tuning, since base models exhibit the same rigidity.
\end{itemize}

\subsection{When to Use Steering vs.\ Prompting}

\begin{table}[H]
\centering
\caption{Updated decision matrix incorporating mid-layer and base model findings.}
\label{tab:decision}
\begin{tabular}{lcc}
\toprule
\textbf{Scenario} & \textbf{Steering} & \textbf{Prompting} \\
\midrule
Base model, mid-layers & \checkmark\checkmark & \checkmark \\
Instruct model, covered task & \checkmark (mid-layer) & \checkmark\checkmark \\
Instruct model, novel task & \checkmark\checkmark & \checkmark \\
Latency-critical (no few-shot budget) & \checkmark\checkmark & $\times$ \\
Compact structured output & $\times$ & \checkmark\checkmark \\
\bottomrule
\end{tabular}
\end{table}

Mid-layer steering is effective even on instruct models, but for tasks within the fine-tuning distribution (calendar extraction), prompt engineering remains simpler and safer. The key risk of steering is the \textbf{narrow coefficient window}: on base models, the gap between ``effective'' ($\alpha=30$, score 0.97) and ``degenerate'' ($\alpha=60$, score 0.33) is a single step.

\subsection{The Tokenizer Bottleneck}

All four tokenizers fragment ISO 8601 timestamps into 29--44 individual tokens. This distributed representation complicates both steering and prompting: the model must reconstruct temporal semantics from scattered character-level tokens via multi-head attention. \textbf{Temporal-aware tokenization} remains an open opportunity.

\subsection{Budget Guidance: A Negative Result}

The null result of budget guidance is informative: instruct models fine-tuned for structured output already produce near-minimal token sequences. Budget guidance is effective when there is \textit{compressible overhead} (e.g., verbose chain-of-thought on reasoning tasks). For agentic tasks with compact outputs, inference-time token control adds no value---the model is already Pareto-optimal.

% ===================================================================
\section{Conclusion}

We presented a comprehensive study of activation steering vs.\ prompt engineering for agentic calendar tasks, going beyond single-layer analysis to reveal the full picture.

\textbf{Mid-layer steering works.} Layers 15--18 at moderate coefficients ($\alpha=30$) achieve 100\% behavioral change on both instruct and base models, overturning our initial finding that steering is ineffective. The key insight is that \textit{where} you steer matters more than \textit{how much}.

\textbf{Late-layer rigidity is architectural.} Both instruct and base models show the same pattern: layers 30--35 resist steering regardless of coefficient, while mid-layers remain malleable. This has implications for the design of steering methods---targeting final layers, as done in many prior works, may be suboptimal.

\textbf{Base models are more steerable but fragile.} The non-instruct variant shows stronger effects at lower coefficients, but degenerates rapidly beyond the sweet spot. Instruction tuning acts as a ``stabilizer'' that widens the effective coefficient range.

\textbf{Budget guidance is task-dependent.} For structured extraction where the instruct model already generates compact outputs, token-level length control adds no value. This is a useful negative result for practitioners.

All code, data, and steering vectors are available at \url{https://github.com/Shumatsurontek/steering-research}.

% ===================================================================
\section*{Acknowledgments}
Analysis performed using the Transformer Explainer tool~\cite{cho2024transformer} and Neuronpedia interpretability platform. Model weights courtesy of the Qwen team~\cite{qwen3}.

% ===================================================================
\bibliographystyle{plain}
\begin{thebibliography}{99}

\bibitem{turner2023activation}
A.~Turner, L.~Thiergart, D.~Udell, G.~Leech, U.~Mini, and M.~MacDiarmid.
\newblock Activation addition: Steering language models without optimization.
\newblock \textit{arXiv preprint arXiv:2308.10248}, 2023.

\bibitem{gao2025steering}
L.~Gao et al.
\newblock Steering {LLM} reasoning through bias-only adaptation.
\newblock In \textit{EMNLP}, 2025.
\newblock arXiv:2505.18706.

\bibitem{li2025budget}
J.~Li, W.~Zhao, Y.~Zhang, and C.~Gan.
\newblock Steering {LLM} thinking with budget guidance.
\newblock \textit{arXiv preprint arXiv:2506.13752}, 2025.

\bibitem{cunningham2023sparse}
H.~Cunningham, A.~Ewart, L.~Riggs, R.~Huben, and L.~Sharkey.
\newblock Sparse autoencoders find highly interpretable features in language models.
\newblock \textit{arXiv preprint arXiv:2309.08600}, 2023.

\bibitem{lieberum2024gemma}
T.~Lieberum et al.
\newblock Gemma {S}cope: Open sparse autoencoders everywhere all at once on {G}emma 2.
\newblock \textit{arXiv preprint arXiv:2408.05147}, 2024.

\bibitem{wei2022chain}
J.~Wei et al.
\newblock Chain-of-thought prompting elicits reasoning in large language models.
\newblock In \textit{NeurIPS}, 2022.

\bibitem{schick2023toolformer}
T.~Schick et al.
\newblock Toolformer: Language models can teach themselves to use tools.
\newblock In \textit{NeurIPS}, 2023.

\bibitem{li2024inference}
K.~Li, O.~Patel, F.~Vi\'{e}gas, H.~Pfister, and M.~Wattenberg.
\newblock Inference-time intervention: Eliciting truthful answers from a language model.
\newblock In \textit{NeurIPS}, 2024.

\bibitem{zou2023representation}
A.~Zou, L.~Phan, S.~Chen, J.~Campbell, P.~Guo, R.~Ren, A.~Pan, X.~Yin, M.~Mazeika, A.~Dombrowski, S.~Goel, N.~Li, M.~Lin, S.~Wang, C.~Xiao, J.~Zou, Z.~Kolter, and D.~Hendrycks.
\newblock Representation engineering: A top-down approach to {AI} transparency.
\newblock \textit{arXiv preprint arXiv:2310.01405}, 2023.

\bibitem{cho2024transformer}
H.~Cho et al.
\newblock Transformer {E}xplainer: Interactive learning of text generative models.
\newblock 2024.
\newblock \url{https://poloclub.github.io/transformer-explainer/}

\bibitem{qwen3}
Qwen Team.
\newblock Qwen3 technical report.
\newblock \textit{arXiv preprint arXiv:2505.09388}, 2025.

\end{thebibliography}

\end{document}
